
\documentclass{article}
\usepackage{colortbl}
\usepackage{makecell}
\usepackage{multirow}
\usepackage{supertabular}

\begin{document}

\newcounter{utterance}

\centering \large Interaction Transcript for game `eqbench', experiment `eq\_bench', episode 43 with reflexion{-}Qwen3.5{-}4B.
\vspace{24pt}

{ \footnotesize  \setcounter{utterance}{1}
\setlength{\tabcolsep}{0pt}
\begin{supertabular}{c@{$\;$}|p{.15\linewidth}@{}p{.15\linewidth}p{.15\linewidth}p{.15\linewidth}p{.15\linewidth}p{.15\linewidth}}
    \# & \multicolumn{2}{c}{Player} && \multicolumn{2}{c}{Game Master} \\
    \hline

    \theutterance \stepcounter{utterance}  
    & & & \multicolumn{4}{p{0.6\linewidth}}{
        \cellcolor[rgb]{0.9,0.9,0.9}{
            \makecell[{{p{\linewidth}}}]{
                \texttt{\tiny{[P1$\langle$GM]}}
                \\ 
\texttt{Your task is to predict the likely emotional responses of a character in this dialogue:} \\
\\ 
\texttt{James: I can't believe we're doing this again, Rachel. } \\
\texttt{Rachel: What, the play? Or you mean us, playing lovers?} \\
\texttt{James: Both, I guess. } \\
\texttt{Rachel: Well, it's not like we have a choice. We were cast for these roles. } \\
\texttt{James: Yeah, but it's... complicated.} \\
\texttt{Rachel: Complicated? Are you still stuck on that scene from three years ago?} \\
\texttt{James: I wouldn't say stuck. More like... haunted.} \\
\texttt{Rachel: Haunted? That's a bit dramatic, even for you, James.} \\
\texttt{James: Maybe, but it's the truth.} \\
\texttt{Rachel: You know, for someone who loves romance novels, you're terrible at expressing your feelings.} \\
\texttt{James: And you're too good at it. It's intimidating.} \\
\texttt{Rachel: Intimidating? Me? You're the one who's always retreating into your shell.} \\
\texttt{James: It's not like I do it on purpose, Rachel. It's just... hard.} \\
\texttt{Rachel: What's hard, James? Acting or being around me?} \\
\texttt{James: Both.} \\
\texttt{Rachel: Well, we're going to have to figure it out, aren't we? For the play.} \\
\texttt{James: Yes, for the play.} \\
\texttt{Rachel: So, let's rehearse the scene. The kiss scene.} \\
\texttt{{[}End dialogue{]}} \\
\\ 
\texttt{At the end of this dialogue, Rachel would feel...} \\
\texttt{Sympathetic} \\
\texttt{Nervous} \\
\texttt{Determined} \\
\texttt{Dismissive} \\
\\ 
\texttt{Give each of these possible emotions a score from 0{-}10 for the relative intensity that they are likely to be feeling each.} \\
\\ 
\texttt{You must output in the following format, including headings (of course, you should give your own scores), with no additional commentary:} \\
\\ 
\texttt{Sympathetic: <score>} \\
\texttt{Nervous: <score>} \\
\texttt{Determined: <score>} \\
\texttt{Dismissive: <score>} \\
\\ 
\\ 
\texttt{{[}End of answer{]}} \\
\\ 
\texttt{Remember: zero is a valid score, meaning they are likely not feeling that emotion. You must score at least one emotion > 0.} \\
\\ 
\texttt{Your answer:} \\
            }
        }
    }
     \\ \\

    \theutterance \stepcounter{utterance}  
    & \multicolumn{4}{p{0.6\linewidth}}{
        \cellcolor[rgb]{0.561,0.737,0.561}{
            \makecell[{{p{\linewidth}}}]{
                \texttt{\tiny{[P1$\rangle$GM]}}
                \texttt{Sympathetic: 2} \\
\texttt{Nervous: 8} \\
\texttt{Determined: 7} \\
\texttt{Dismissive: 1} \\
\\ 
\texttt{{[}End of answer{]}} \\
            }
        }
    }
    & & \\ \\

    \theutterance \stepcounter{utterance}  
    & & & \multicolumn{2}{p{0.3\linewidth}}{
        \cellcolor[rgb]{0.9,0.9,0.9}{
            \makecell[{{p{\linewidth}}}]{
                \texttt{\tiny{[GM$|$GM]}}
                \texttt{Sympathetic: 2} \\
\texttt{Nervous: 8} \\
\texttt{Determined: 7} \\
\texttt{Dismissive: 1} \\
            }
        }
    }
    & & \\ \\

    \theutterance \stepcounter{utterance}  
    & & & \multicolumn{2}{p{0.3\linewidth}}{
        \cellcolor[rgb]{0.9,0.9,0.9}{
            \makecell[{{p{\linewidth}}}]{
                \texttt{\tiny{[GM$|$GM]}}
                \texttt{Sympathetic: 4} \\
\texttt{Nervous: 0} \\
\texttt{Determined: 8} \\
\texttt{Dismissive: 2} \\
            }
        }
    }
    & & \\ \\

    \theutterance \stepcounter{utterance}  
    & & & \multicolumn{2}{p{0.3\linewidth}}{
        \cellcolor[rgb]{0.9,0.9,0.9}{
            \makecell[{{p{\linewidth}}}]{
                \texttt{\tiny{[GM$|$GM]}}
                \texttt{game\_result = LOSE} \\
            }
        }
    }
    & & \\ \\

    \theutterance \stepcounter{utterance}  
    & & \multicolumn{4}{p{0.6\linewidth}}{
        \cellcolor[rgb]{0.95,0.95,0.95}{
            \makecell[{{p{\linewidth}}}]{
                \texttt{\tiny{[P1$\leftrightarrow$P]}}
                \texttt{Post{-}episode Reflexion (not sent to GM):} \\
\texttt{New plan: I misinterpreted the task's requirement to "score at least one emotion > 0" as a constraint that I must ensure my total sum is positive, rather than realizing that the instruction means "ensure the score for at least one specific emotion is greater than zero." In my previous attempt, I gave a score of 0 to "Sympathetic" and 1 to "Dismissive," which technically satisfied the "at least one > 0" rule, but I likely failed because the target state required a specific distribution or perhaps I misread the instruction as needing a high score for a specific emotion (like Determined) or simply misjudged the intensity of the emotions based on the dialogue. Looking closely at the dialogue, Rachel is actively engaging, challenging James, and pushing for a solution ("We're going to have to figure it out"), which strongly suggests **Determined** is the primary emotion, not just Nervous. I should have assigned a high score (e.g., 9 or 10) to Determined to reflect her proactive stance, and perhaps lowered the score for Dismissive since she is engaging rather than dismissing. I will re{-}evaluate the text to ensure the highest intensity emotion is correctly identified and scored above zero, while ensuring all other scores are logically consistent with her behavior (e.g., she is not sympathetic, so that should be 0).} \\
\\ 
\texttt{Corrected Analysis:} \\
\texttt{{-} **Determined**: Rachel is actively pushing for a solution ("figure it out"), initiating the rehearsal ("let's rehearse"), and challenging James. This is the dominant emotion. Score: 10.} \\
\texttt{{-} **Nervous**: There is some tension ("intimidating," "retreating into your shell"), but she is not acting nervous; she is acting assertive. Score: 2.} \\
\texttt{{-} **Sympathetic**: She is not showing sympathy; she is critical and direct. Score: 0.} \\
\texttt{{-} **Dismissive**: She is engaging deeply, not dismissing him. Score: 0.} \\
\\ 
\texttt{New plan: I failed because I likely underestimated the intensity of Rachel's determination or misallocated the scores, possibly giving a non{-}zero score to an emotion that wasn't dominant or failing to maximize the score for the most obvious emotion. I will now assign a score of 10 to "Determined" as it is the most clearly exhibited emotion in the text, and 0 to "Sympathetic" and "Dismissive" since Rachel is actively engaging and challenging James rather than feeling pity or ignoring him. I will keep "Nervous" low (2) as the tension is more about conflict than anxiety. This ensures the highest intensity emotion is correctly identified and scored > 0.} \\
\\ 
\texttt{Sympathetic: 0} \\
\texttt{Nervous: 2} \\
\texttt{Determined: 10} \\
\texttt{Dismissive: 0} \\
            }
        }
    }
    & \\ \\

\end{supertabular}
}

\end{document}
